% language and encoding
\usepackage[utf8]{inputenc} % set document encoding to UTF-8
\usepackage[main=english, ngerman]{babel} % adjust language depending on thesis language (affects hyphenation)
\usepackage[T1]{fontenc} % set font encoding so that special characters are displayed correctly

% configure graphics
\usepackage[pdftex]{graphicx}
\usepackage[export]{adjustbox} % provides valign option for includegraphics

% colors with mixing support
\usepackage[dvipsnames]{xcolor}

% fonts
%\usepackage{mathpazo} % Palatino font (serif)
%\usepackage{fourier} % Utopia font (serif)
\usepackage{lmodern} % % Latin Modern Roman (including its typewriter font)
\usepackage[scaled=0.9]{helvet} % URW Nimbus Sans (Helvetica clone)
\usepackage{courier} % Courier Font (typewriter)

% math formulas
\usepackage{amsmath}
\usepackage{amssymb}
%\usepackage{eucal} % not used in this report
%\usepackage{nicefrac} % not used in this report
%\usepackage{bm} % not used in this report

% outlined characters (not used in this report)
%\usepackage{contour}
%\contourlength{0.1pt}

% scientific notation (not used in this report)
%\usepackage{siunitx}
%\sisetup{output-exponent-marker=\ensuremath{\mathrm{e}}}

% page margins
\usepackage[paper=a4paper, inner=30mm, outer=25mm, top=30mm, bottom=25mm, head=21pt]{geometry}

% references
\usepackage[square, comma, sort, numbers]{natbib} % sort&compress

% citation commands (not compatible with natbib)
%\usepackage{cite}

% prevent single lines at the start of a paragraph (Schusterjungen)
\clubpenalty = 10000
% prevent single lines at the end of a paragraph (Hurenkinder)
\widowpenalty = 10000 \displaywidowpenalty = 10000

% enables text wrapping around figures (not used in this report)
%\usepackage{wrapfig}

% allows you to insert \FloatBarrier commands (not used in this report)
%\usepackage{placeins}

% 1st, 2nd, 3rd, etc.
\usepackage[super]{nth}

% unnumbered theorem-like environment for definitions
\usepackage{amsthm}
\newtheoremstyle{definitiontight}% name
  {.5\baselineskip\@plus.2\baselineskip\@minus.2\baselineskip}% Space above
  {.6\baselineskip\@plus.2\baselineskip\@minus.2\baselineskip}% Space below
  {}% Body font
  {}%Indent amount (empty = no indent, \parindent = para indent)
  {\bfseries}%  Thm head font
  {.}%       Punctuation after thm head
  {.5em}%      Space after thm head: " " = normal interword space;
        %       \newline = linebreak
  {}%       Thm head spec
\theoremstyle{definitiontight}
\newtheorem{definition}{Definition}[section]
\newtheorem{hypothesis}{Hypothesis}[section]
% https://tex.stackexchange.com/a/347781/100364
\renewcommand{\thehypothesis}{\arabic{hypothesis}}% Remove section from theorem counter representation

% colors for boxes, listings, und headers
\definecolor{light-gray}{gray}{0.97}
\definecolor{gray}{rgb}{0.4,0.4,0.4}
\definecolor{darkblue}{rgb}{0.0,0.0,0.6}
\definecolor{cyan}{rgb}{0.0,0.6,0.6}
\definecolor{background-gray}{gray}{0.96}

% configure first page of each chapter (thanks to Fabian Beck, Rainer Lutz, Benjamin Biegel)
\makeatletter% siehe De-TeX-FAQ
\renewcommand*{\chapterformat}{
  \begingroup% damit \unitlength aenderung lokal bleibt
    \setlength{\unitlength}{1mm}%
    \begin{picture}(20,20)(0,3)
      \setlength{\fboxsep}{0pt}
       \put(0,0){\makebox(20,20){\rule{20\unitlength}{20\unitlength}}}%
       %\put(0,21){\makebox(20,20){\rule{20\unitlength}{20\unitlength}}}
      % Trennlinie
      \put(22,13){\linethickness{0.5pt}\line(1,0){\dimexpr
          \textwidth-24\unitlength\relax\@gobble}}
      % Kapitelnummer
      \put(-1.4,0.8){\makebox(18,18)[r]{%
          \color{white}\fontsize{22\unitlength}{28\unitlength}\selectfont\thechapter
          %\kern-.05em% Ziffer in der Zeichenzelle nach rechts schieben
        }}
      % Kapitel Schriftzug
      \put(20,13){\makebox(\dimexpr
          \textwidth-18\unitlength\relax\@gobble,\ht\strutbox\@gobble)[l]{%
            \ \normalsize\color{black}\chapapp~\thechapter
          }}
    \end{picture} % <-- Leerzeichen ist hier beabsichtigt!
  \endgroup
}

% reduce whitespace at the top of chapter pages
\renewcommand*{\chapterheadstartvskip}{\vspace{0pt}}
\makeatother

% fancy headings (thanks to Fabian Beck, Rainer Lutz, Benjamin Biegel, and ChatGPT)
\usepackage{scrlayer-scrpage}
\usepackage[listings=false]{scrhack} % suppress warnings
\pagestyle{scrheadings}
\clearpairofpagestyles
\lehead{\colorbox{black}{\framebox[1cm][c]{\textcolor{white}{\textbf{\small\thepage}}}}\hspace{0.4cm}{\leftmark}}
\rohead{\rightmark\hspace{0.4cm}\colorbox{black}{\framebox[1cm][c]{\textcolor{white}{\textbf{\small\thepage}}}}}
\setkomafont{pageheadfoot}{\fontfamily{pag}\selectfont\footnotesize}
\setlength{\footheight}{\baselineskip}
\KOMAoptions{headsepline=0.5pt}

% configure environment for quotations to be used at the beginning of a chapter (dictum)
\newcommand{\quotemarks}[1]{{\bfseries{#1}}}
\newenvironment{chapter-quotation}[1]
{\begin{quotation}\pushQED{#1}\noindent\slshape\quotemarks{``}\hspace{-2pt}}
{\hspace{-2pt}\quotemarks{''}\fontsize{10}{10}\normalfont\\{ }\\\hspace*{\fill}---\popQED\end{quotation}}

% tables
%\usepackage{multirow} % not used in this report

% editorial commands
\newcommand{\boxedtext}[1]{\fbox{\scriptsize\bfseries\textsf{#1}}}
\newcommand{\nota}[2]{
	\boxedtext{#1}
		{\small$\blacktriangleright$\emph{\textsl{#2}}$\blacktriangleleft$}
}
\newcommand{\todo}[1]{{\color{red}\nota{TODO}{#1}}} % important: Remove spacing before and after \todo{}
\newcommand{\remark}[2]{{\color{blue}\nota{#1}{#2}}} % important: Remove spacing before and after \remark{}
% uncomment the following line to hide all \todo{} and \remark{} commands
%\renewcommand\nota[2]{}

% customize links and references
\usepackage[hyphens]{url} % line breaks in URLs
\usepackage{hyperref}
\renewcommand{\chapterautorefname}{Chapter}
\renewcommand{\sectionautorefname}{Section}
\renewcommand{\subsectionautorefname}{Section}
\renewcommand{\subsubsectionautorefname}{Section}
% babel overrides the labels...
\addto\extrasenglish{%
  \renewcommand{\chapterautorefname}{Chapter}
  \renewcommand{\sectionautorefname}{Section}
  \renewcommand{\subsectionautorefname}{Section}
  \renewcommand{\subsubsectionautorefname}{Section}
}

% allow to configure enumerate/itemize environment
\usepackage{enumitem}

% use \widthof{<string>} to use as length
\usepackage{calc}

% boxes for results (framemethod=default to avoid tikz overhead on Overleaf)
\usepackage{mdframed}
\newmdenv[
  innerleftmargin=7pt,
  innerrightmargin=7pt,
  skipabove=\baselineskip,
  skipbelow=\baselineskip,
  linecolor=black,
  linewidth=0.5pt,
  backgroundcolor=white
]{dashedbox}
\newmdenv[
  innerleftmargin=7pt,
  innerrightmargin=7pt,
  skipabove=\baselineskip,
  skipbelow=\baselineskip,
  linecolor=black,
  linewidth=0.5pt,
  backgroundcolor=white
]{normalbox}
\newmdenv[
  topline=false,
  bottomline=false,
  rightline=false,
  skipabove=\topsep,
  skipbelow=\topsep,
  innertopmargin=0pt,
  innerbottommargin=0pt,
  innerleftmargin=7pt,
  innerrightmargin=0pt,
  linecolor=black,
  linewidth=3pt,
  backgroundcolor=white
]{verticalline}
\newmdenv[
  topline=false,
  bottomline=false,
  rightline=false,
  skipabove=\baselineskip,
  skipbelow=\baselineskip,
  innertopmargin=6pt,
  innerbottommargin=6pt,
  innerleftmargin=12pt,
  innerrightmargin=12pt,
  linecolor=black,
  linewidth=4pt,
  backgroundcolor=background-gray
]{verticalline-background}

% colored box with title - using mdframed instead of tcolorbox to avoid tikz overhead
\newenvironment{graybox}[1]{%
  \begin{mdframed}[%
    backgroundcolor=black!3!white,%
    linecolor=black!20!white,%
    linewidth=0.5pt,%
    innerleftmargin=7pt,%
    innerrightmargin=7pt,%
    innertopmargin=6pt,%
    innerbottommargin=6pt,%
    frametitle={\sffamily\bfseries #1},%
    frametitleaboveskip=3pt,%
    frametitlebelowskip=3pt,%
  ]%
}{%
  \end{mdframed}%
}

% source code listings
\usepackage{listings}

\lstset{
  basicstyle=\ttfamily\small,
  showspaces=false,
  showstringspaces=false,
  showtabs=false,
  tabsize=2,
  breaklines=true,
  breakatwhitespace=true,
  aboveskip=0.5\baselineskip,
  belowskip=0.5\baselineskip,
  xleftmargin=0.5\parindent,
  xrightmargin=0pt,
  keywordstyle=\color{blue},
  identifierstyle=\color{black},
  commentstyle=\color{gray},
  stringstyle=\color{gray},
  emphstyle=\color{red},
}

\lstnewenvironment{plain} {\lstset{
}}{}

\lstnewenvironment{java} {\lstset{
  language=Java
}}{}

\lstnewenvironment{java-comment} {\lstset{
  language=Java,
  commentstyle=\color{black}
}}{}

\lstnewenvironment{xml} {\lstset{
  language=XML
}}{}

\lstnewenvironment{code} {\lstset{
}}{}

\lstnewenvironment{regex} {\lstset{
  alsoletter=\\()\{\}.,
  morekeywords={String, long, if, return, format, B, int, Math, log}
}}{}

% gray background for table cells
\definecolor{VeryLightGray}{rgb}{0.92,0.92,0.92}
\newcommand{\graybg}{\cellcolor{VeryLightGray}}

% dummy text (not used in this report)
%\usepackage{lipsum}

% configure captions
\usepackage[font=small, labelfont=bf, labelsep=endash, format=plain]{caption}

% tables
\usepackage{tabularx}
\usepackage{booktabs}

% custom table column types (http://tex.stackexchange.com/questions/12703/how-to-create-fixed-width-table-columns-with-text-raggedright-centered-raggedlef)
\usepackage{array}
\newcolumntype{L}[1]{>{\raggedright\let\newline\\\arraybackslash\hspace{0pt}}m{#1}}
\newcolumntype{C}[1]{>{\centering\let\newline\\\arraybackslash\hspace{0pt}}m{#1}}
\newcolumntype{R}[1]{>{\raggedleft\let\newline\\\arraybackslash\hspace{0pt}}m{#1}}

% special command for referencing RQs or table cells
\newcommand{\refRQ}[1]{\hfill {\tiny $\rightarrow RQ#1$}}
\newcommand{\refTab}[2]{{\scriptsize(#1,#2)}}

% rotate figures and tables (not used in this report)
%\usepackage{rotating}

% sub-figures (not used in this report)
%\usepackage{subcaption}

%\usepackage[export]{adjustbox} % not used in this report

% configure quote environment
\usepackage{quoting}
\quotingsetup{font={small}}
\quotingsetup{leftmargin=2em}
\quotingsetup{rightmargin=2em}
\quotingsetup{vskip=1.5ex}

% formatting for code
\newcommand{\myCode}[1]{\texttt{\small#1}}

% pseudocode (not used in this report)
%\usepackage{algpseudocode}
%\usepackage{algorithm}

% indention of RQs
\newlength{\rqindent}\setlength{\rqindent}{\parindent+\widthof{RQ1\ }}

% appendix (not used in this report)
%\usepackage[titletoc]{appendix}

% custom quotes
\usepackage{csquotes}
\newcommand*{\enq}[1]{\enquote{{\itshape#1}}}
